\chapter{总体研究方案与验证体系}\label{chap:middle-design}
\section{存量平台上的功能分工}

原 AUV 已有以 PC104 和 VxWorks 为核心的底层控制与设备接口。新增 Jetson 计算平台用于承载需要较高计算资源的感知、融合、规划与优化控制，上位机承担任务配置、状态观察和实验记录。这样的分工保留执行侧既有接口，允许新增算法以受约束的设定值进入执行链，并使任务决策与设备反馈能够分别审查。本文所研究的双脑结构由此服务于存量平台的逐步改造，而不是要求在一次集成中替换全部控制软件。

图\ref{fig:middle-architecture}概括了数据和命令之间的关系。图中表示的是系统设计与验证接口；各条链路达到的实验层级在第\ref{ch:middle-systems}章分别说明。磁采集、声学线索和导航信息进入感知与估计侧，决策模块根据质量和任务状态选择搜索、横切确认或沿线跟踪；设定值经控制和仲裁后才进入执行脑。异常状态沿反馈链向上报告，同时参与限制后续命令。

\begin{figure}[htbp]
\centering
\begin{tikzpicture}[font=\small,box/.style={draw,rounded corners=2pt,align=center,
  text width=3.7cm,minimum height=1.05cm},arr/.style={-{Latex[length=2mm]},thick}]
\node[box] (sensors) {磁采集、导航与声学线索\\观测、时间戳、质量};
\node[box,right=1.0cm of sensors] (estimate) {艇载感知与状态估计\\几何、状态、协方差};
\node[box,below=0.7cm of estimate] (decision) {任务决策与控制\\阶段选择、运动设定值};
\node[box,left=1.0cm of decision] (execute) {仲裁、PC104 与执行机构\\命令准入、反馈与异常};
\node[box,below=0.8cm of execute,text width=8.5cm,xshift=2.35cm] (record) {上位机与记录链：任务配置、运行状态、原始数据及回放};
\draw[arr] (sensors)--(estimate);
\draw[arr] (estimate)--(decision);
\draw[arr] (decision)--(execute);
\draw[arr] (execute)--(sensors);
\draw[arr] (execute)--(execute.south|-record.north);
\draw[arr] (decision)--(decision.south|-record.north);
\end{tikzpicture}
\caption{感知、决策与执行的接口关系示意；示意连接不表示整艇联调已经完成}
\label{fig:middle-architecture}
\end{figure}

改造工作还受到物理布局的约束。推进器、电源转换和计算设备可能在磁传感器附近产生干扰，载荷安装又会改变观测点与导航参考点之间的几何关系。因此，远离主要干扰源、固定传感器外参和保持可重复布线，是测量方案的一部分。原论文提出的舾装与磁卫生设计为集成提供依据；尚缺相同工况屏蔽前后测量的设计项不作为已经获得定量收益的结果。

\section{多源观测与状态信息的衔接}

系统采用局部导航坐标描述航行器运动，惯性测量提供短时传播，多普勒测速仪（Doppler velocity log，DVL）、深度及其他外部观测对漂移进行约束。磁传感器固定于艇体，其测量位置与导航参考点存在杆臂，方向也由安装姿态决定。设导航系为 $n$、艇体系为 $b$，导航参考点位置为 $\boldsymbol p_b^n$，磁传感器杆臂为 $\boldsymbol l_m^b$，则观测点位置满足
\begin{equation}
  \boldsymbol p_m^n=\boldsymbol p_b^n+\boldsymbol R_b^n\boldsymbol l_m^b.
  \label{eq:middle-lever}
\end{equation}
式中 $\boldsymbol R_b^n$ 将艇体系向量旋转至导航系。主动横切使用的是磁观测点的空间坐标；若把艇体参考点直接当作探头位置，姿态变化便会转换成剖面位置误差。测量链标定、运动记录与算法几何必须使用一致的参考关系。

声学输入在本课题中主要提供目标方向、横向位置或候选区域。侧扫结果具有斜距和海床投影关系，其输出需要连同来源、时间与质量解释。视觉标定则用于实验室中建立传感器与导线的相对几何，不等同于 AUV 水下任务的在线真值来源。将这些角色分开后，可以明确哪些信息用于算法输入，哪些只用于结果评价，以及哪些仍需由独立调查建立目标身份。

状态估计不仅传递位置和姿态，还需要向控制层传递不确定性。观测质量下降可能通过增大观测噪声、鲁棒更新或拒绝异常值影响滤波；控制侧据此调节跟踪与控制代价，并通过速率与带宽约束限制指令变化。但质量评分、滤波协方差与任务成功率并非同一量，不能把一个归一化评分直接当成真实可靠性概率。第\ref{ch:middle-systems}章以对照结果说明这种信息利用的收益与代价。

\section{巡检流程与结果发布原则}

磁巡检 v2 以先验区域搜索为入口，经历磁检测、主动横切确认、显式 FWHM 估计和沿缆跟踪。搜索使用可调宽度、测线间距和速度在给定边界内生成航点，目的是让初始位置误差转化为可规划的覆盖范围。检测通过滑动窗口、阈值和滞回判断是否存在值得横切的磁异常，确认阶段再主动建立空间剖面。这样的顺序将“可能发现异常”和“已有几何证据”分开，减少单次大幅值触发直接进入跟踪的风险。

横切轨迹需要兼顾观测与运动约束。扫描过窄时，峰值两侧可能都不足以达到半高点；扫描幅度扩大可以增加空间支撑，却同时增加航程和转向需求。对于有最小转弯半径与低速响应限制的 AUV，理想折线也不能直接视作可执行航迹。因此本阶段将磁几何可辨识、横切完成、接入沿线和有限任务到端分别检查，并在仿真中保留实际运动产生的回归段。

结果发布遵循分量独立的原则。磁剖面的中心可以提供横向信息，场形系数和离底高度则共同决定能否发布埋深；目标身份还需要与独立先验或声学线索关联。若只有可解释的横向中心而缺少可靠深度条件，系统保留横向信息，将深度记为空值。若目标关联或导航条件不满足指定任务要求，则不以“检测到磁异常”替代“已完成指定海缆巡检”。

执行许可同样需要独立判断。传感器输入、估计输出、任务授权和控制命令可能采用不同的发布频率，故障后的零设定值必须传播到后续节点并被实际接收。另一方面，即使各层命令都归零，航行器仍可能因惯性或水流继续运动。本文将命令撤回、任务中断与物理停止分开记录，避免由软件状态推断未测量的运动结果。

\section{从可重复仿真到实物的验证层次}

本研究采用逐层验证而非单一综合分数。离线磁剖面对照用于隔离信号、增益和场形因素；运动仿真用于检验搜索、转向、接入与跟踪之间的耦合；软件链路与独立记录用于确认消息和故障的传播；磁采集、艇载求解与真实 PC104 通信则提供物理接口锚点。表\ref{tab:middle-evidence-layers}列出这些层次所支持的主张及其后续衔接。

\begin{table}[htbp]
\centering\caption{分层验证的研究问题与输出}\label{tab:middle-evidence-layers}
\begin{tabularx}{\linewidth}{p{2.3cm}XX}\toprule
验证层次&本阶段主要问题&用于后续工作的输出\\\midrule
离线算法对照&场形、噪声与增益怎样改变反演&误差、有效率、拒绝原因及适用条件\\
闭环运动仿真&观测轨迹能否实际形成，何时可以跟踪&运动航迹、真实横偏、任务阶段与退出状态\\
软件及协议链&正常与故障信息能否按接口传播&同次状态、命令、接收与记录证据\\
实物子链&采集、计算与执行脑接口怎样运行&噪声和拟合、求解时序、通信与故障响应\\
整艇及任务试验&装配后的载体与任务链是否协同有效&受控运动、完整巡检、恢复处置及重复性\\\bottomrule
\end{tabularx}
\end{table}

动力学仿真采用仓库已有的 Python Vehicle Simulator（PVS）路线，运动模型与控制接口围绕海洋航行器建模方法组织\cite{fossen2011handbook}。不同实验中的代理模型、原生六自由度回路和 ROS 通信链分别保留其版本，不把多个时段的结果拼接为同一次端到端验收。当前报告优先选取能够回答核心问题的冻结结果，对源码快照完整性较弱或原始记录不在当前检出的材料，内部证据索引记录其支撑程度。
